\section{Руководство пользователя}

\subsection{Вход в систему}

Доступ к приложению осуществляется через веб-браузер по адресу, заданному при развёртывании системы.
При первом обращении пользователь перенаправляется на страницу входа, на которой расположены поля ввода имени пользователя и пароля, переключатель языка интерфейса (русский или английский) и кнопка <<Войти>>.
Внешний вид страницы входа представлен на рисунке~\ref{fig:login}.

\begin{figure}[htbp]
\centering
\figureorstub{login}{скриншот страницы входа}
\caption{Страница входа в систему}
\label{fig:login}
\end{figure}

При успешной аутентификации пользователь перенаправляется на главную страницу приложения~--~страницу результатов отбора.
При неверных учётных данных на странице отображается сообщение об ошибке; каждая попытка входа фиксируется в журнале авторизации независимо от её исхода.
Форма проверяет заполненность обоих полей на стороне клиента: при попытке отправить пустую форму поля подсвечиваются с подсказкой об обязательности заполнения, а запрос на сервер не отправляется.
Переключатель языка над формой позволяет выбрать язык интерфейса до входа в систему; выбор сохраняется в локальном хранилище браузера и применяется автоматически при последующих сеансах.
Выбранный язык влияет не только на отображение текстов клиентской части: в заголовок каждого HTTP-запроса подставляется значение \code{Accept-Language}, благодаря чему сообщения об ошибках с сервера также возвращаются на выбранном языке.
После входа в систему пользователь может в любой момент сменить пароль через меню учётной записи в верхней части экрана.

\subsection{Главное меню и навигация}

Главное меню расположено в левой боковой панели приложения и содержит следующие пункты: <<Результаты>>, <<Абитуриенты>>, <<Факультеты>>, <<Кафедры>>, <<Специальности>>, <<Критерии оценивания>>, <<Группы критериев>>, <<Аудит>>, <<Журнал аудита>>, <<Пользователи>>.
Видимость пунктов меню одинакова для всех аутентифицированных пользователей; ограничения по ролям применяются на уровне конкретных операций (создание, изменение, удаление).
Внешний вид главного меню можно наблюдать в левой боковой панели на рисунке~\ref{fig:faculties}.

В верхней части каждой страницы расположен заголовок с переключателем темы (светлая или тёмная), переключателем языка интерфейса (русский или английский) и меню текущего пользователя.
Меню пользователя содержит пункты <<Сменить пароль>> и <<Выход>>.
Пункт <<Сменить пароль>> открывает диалоговое окно смены пароля, в котором необходимо ввести текущий пароль, новый пароль и его повтор.
Пункт <<Выход>> удаляет токен доступа из локального хранилища браузера и возвращает пользователя на страницу входа.

Переключение языка интерфейса применяется немедленно без перезагрузки страницы.
Одновременно с переводом текстов клиентской части в заголовок \code{Accept-Language} последующих HTTP-запросов подставляется выбранный язык, благодаря чему сервер возвращает сообщения об ошибках на том же языке.

\subsection{Работа со структурой университета}

Страницы <<Факультеты>>, <<Кафедры>> и <<Специальности>> построены по единому шаблону: таблица с серверной пагинацией и фильтрами в верхней части, кнопка <<Добавить>> над таблицей и кнопки <<Изменить>>, <<Удалить>> в каждой строке.
Кнопки записи отображаются только пользователям с правом \code{WriteStructure} (роли \code{SuperAdmin} и \code{DataAdministrator}).
Внешний вид страницы факультетов представлен на рисунке~\ref{fig:faculties}.

\begin{figure}[htbp]
\centering
\figureorstub{faculties}{скриншот страницы факультетов}
\caption{Страница факультетов}
\label{fig:faculties}
\end{figure}

Для добавления факультета пользователь нажимает кнопку <<Добавить>>, заполняет поля наименования и сокращённого наименования и подтверждает создание.
Аналогично создаются кафедры (с обязательным указанием факультета) и специальности (с указанием кафедры).
При попытке удалить факультет с привязанными кафедрами или кафедру с привязанными специальностями отображается сообщение об ошибке~--~целостность ссылок гарантируется внешними ключами на уровне базы данных.

Страница редактирования специальности содержит дополнительные вкладки с привязанными к специальности конкурсными списками и категориями приёма.
Для каждого конкурсного списка указывается его наименование (например, <<Бюджет>>, <<Платное>>) и общий план приёма.
В пределах конкурсного списка администратор формирует категории приёма: для каждой категории указывает её наименование, квоту мест, приоритет и группу критериев оценивания.
Внешний вид страницы редактирования специальности представлен на рисунке~\ref{fig:specialtyedit}.

\begin{figure}[htbp]
\centering
\figureorstub{specialty-edit}{скриншот страницы редактирования специальности}
\caption{Страница редактирования специальности}
\label{fig:specialtyedit}
\end{figure}

\subsection{Работа с критериями оценивания}

Страница <<Критерии оценивания>> содержит перечень всех показателей, по которым оцениваются абитуриенты.
Для каждого критерия указывается наименование, диапазон допустимых значений (поля \code{minvalue} и \code{maxvalue}) и тип: <<больше~--~лучше>> (\code{higher\_is\_better}) или <<меньше~--~лучше>> (\code{lower\_is\_better}).
Внешний вид страницы критериев оценивания представлен на рисунке~\ref{fig:criteria}.

\begin{figure}[htbp]
\centering
\figureorstub{evaluation-criteria}{скриншот страницы критериев оценивания}
\caption{Страница критериев оценивания}
\label{fig:criteria}
\end{figure}

Страница <<Группы критериев>> позволяет объединять критерии в именованные наборы, привязываемые впоследствии к категориям приёма.
В составе группы каждый критерий получает значение поля \code{priority}, определяющего порядок лексикографической сортировки кандидатов: критерий с меньшим значением \code{priority} учитывается в сравнении первым.
Внешний вид страницы редактирования группы критериев представлен на рисунке~\ref{fig:criteriagroup}.

\begin{figure}[htbp]
\centering
\figureorstub{evaluation-criteria-group-edit}{скриншот страницы редактирования группы критериев}
\caption{Страница редактирования группы критериев}
\label{fig:criteriagroup}
\end{figure}

\subsection{Работа с абитуриентами}

Страница <<Абитуриенты>> содержит таблицу всех зарегистрированных абитуриентов с возможностью фильтрации по имени, внешнему идентификатору документа и статусу валидации.
Внешний вид страницы представлен на рисунке~\ref{fig:applicants}.

\begin{figure}[htbp]
\centering
\figureorstub{applicants}{скриншот страницы списка абитуриентов}
\caption{Страница списка абитуриентов}
\label{fig:applicants}
\end{figure}

Для добавления нового абитуриента пользователь с правом \code{WriteApplicants} нажимает кнопку <<Добавить>> и заполняет форму: имя, внешний идентификатор документа, текстовые заметки.
Если в системе уже существует активная или удалённая запись с тем же внешним идентификатором, отображается сообщение об ошибке~--~уникальность внешнего идентификатора обеспечивается ограничением базы данных.

После сохранения абитуриента открывается страница его редактирования, на которой расположены вкладки для ввода оценок по критериям и для управления заявками в категории приёма.
Внешний вид страницы редактирования абитуриента представлен на рисунке~\ref{fig:applicantedit}.

\begin{figure}[htbp]
\centering
\figureorstub{applicant-edit}{скриншот страницы редактирования абитуриента}
\caption{Страница редактирования абитуриента}
\label{fig:applicantedit}
\end{figure}

На вкладке оценок отображается список всех критериев, доступных в системе.
Для каждого критерия пользователь вводит численное значение, которое автоматически проверяется на принадлежность диапазону, заданному для критерия.
При вводе значения вне диапазона выводится сообщение об ошибке с указанием границ.
При сохранении формы выполняется отправка изменений на сервер, после чего автоматически запускается пересчёт результатов отбора.

На вкладке заявок пользователь добавляет записи о подаче заявления в конкретные категории приёма с указанием приоритета предпочтения.
Один абитуриент может подать заявки в несколько категорий разных специальностей.
Поле приоритета предпочтения принимает целые положительные значения; меньшее значение соответствует более предпочтительной для абитуриента категории.

Удаление абитуриента выполняется в два этапа.
При нажатии кнопки <<Удалить>> в строке абитуриента создаётся запрос на удаление со статусом \code{pending}; абитуриент немедленно скрывается из всех таблиц и из результатов отбора, но физически из базы данных не удаляется.
Окончательное удаление либо отмена запроса выполняются на странице <<Аудит>> пользователем с правом \code{WriteAudit}.

\subsection{Просмотр результатов отбора}

Страница <<Результаты>> содержит список всех конкурсных списков системы с указанием специальности, общего плана приёма и количества зачисленных абитуриентов.
Внешний вид страницы представлен на рисунке~\ref{fig:results}.

\begin{figure}[htbp]
\centering
\figureorstub{results}{скриншот страницы результатов отбора}
\caption{Список результатов отбора}
\label{fig:results}
\end{figure}

Для каждой строки списка доступно скачивание результатов соответствующего конкурсного списка в формате Excel через ссылку <<Скачать Excel>>.
Сформированный файл содержит несколько листов~--~один лист на категорию приёма~--~и включает столбцы с именем абитуриента, внешним идентификатором документа, приоритетом предпочтения заявки и значениями всех критериев из группы категории.
Имя файла формируется автоматически на основе наименования специальности и конкурсного списка.

При выборе конкретного конкурсного списка открывается страница детального просмотра с разбиением по категориям.
Для каждой категории отображается её квота, приоритет защиты и список зачисленных абитуриентов с их оценками по критериям.
Внешний вид страницы детального просмотра представлен на рисунке~\ref{fig:competitionresult}.

\begin{figure}[htbp]
\centering
\figureorstub{competition-list-result}{скриншот страницы детального просмотра результатов конкурсного списка}
\caption{Детальный просмотр результатов конкурсного списка}
\label{fig:competitionresult}
\end{figure}

Кнопка <<Пересчитать>> над списком конкурсных списков позволяет вручную запустить пересчёт результатов отбора.
Кнопка доступна пользователям с правом \code{WriteStructure}.
В штатном режиме ручной запуск не требуется: пересчёт выполняется автоматически при каждом изменении данных абитуриентов, их оценок, заявок и при изменении состава критериев в группах.
Ручной запуск используется в случае разовых административных изменений данных непосредственно в базе данных в обход интерфейса, а также для пересчёта по запросу при отладке.

\subsection{Подсистема аудита}

Страница <<Аудит>> содержит пять вкладок, каждая из которых отображает один из видов оповещений системы.
Последовательно просматривая вкладки, оператор получает полную картину состояния данных, требующих проверки или обработки.
Внешний вид страницы аудита представлен на рисунке~\ref{fig:audit}.

\begin{figure}[htbp]
\centering
\figureorstub{audit}{скриншот страницы аудита}
\caption{Страница аудита}
\label{fig:audit}
\end{figure}

Вкладка <<Невалидированные абитуриенты>> отображает список абитуриентов, статус валидации которых равен \code{false}.
Пользователь с правом \code{WriteAudit} может подтвердить валидацию (\code{validated = true}) или вновь снять её для каждого абитуриента.
При последующем изменении любого поля данных абитуриента (имени, внешнего идентификатора, оценок, заявок) статус автоматически возвращается в значение \code{false}, и абитуриент вновь появляется в этой вкладке.

Вкладка <<Неполные данные>> отображает абитуриентов, у которых не заполнены оценки по всем критериям, требуемым категориями, в которые они подали заявки.
Для каждого такого абитуриента показывается список отсутствующих критериев~--~это позволяет оператору быстро устранить пробелы в данных.

Вкладка <<Некорректные группы критериев>> отображает группы критериев, в которых отсутствуют обязательные элементы или указаны неверные приоритеты.

Вкладка <<Некорректные конкурсные списки>> отображает конкурсные списки, в которых обнаружены проблемы с категориями приёма: отсутствие группы критериев, недопустимые значения квот или приоритетов.

Вкладка <<Запросы на удаление>> отображает запросы со статусом \code{pending}.
Для каждого запроса доступны кнопки <<Подтвердить>> и <<Отклонить>>.
При подтверждении абитуриент скрывается окончательно; при отклонении абитуриент возвращается в активные данные, помечается как невалидированный и автоматически запускается пересчёт результатов отбора.
После отклонения возможно создание нового запроса на удаление того же абитуриента.
Внешний вид вкладки запросов на удаление представлен на рисунке~\ref{fig:pendingdeletions}.

\begin{figure}[htbp]
\centering
\figureorstub{pending-deletions}{скриншот вкладки запросов на удаление}
\caption{Вкладка запросов на удаление страницы аудита}
\label{fig:pendingdeletions}
\end{figure}

Страница <<Журнал аудита>> отображает полный журнал всех действий пользователей с возможностью фильтрации по дате, действию, типу сущности и пользователю.
Внешний вид страницы журнала аудита представлен на рисунке~\ref{fig:auditlog}.

\begin{figure}[htbp]
\centering
\figureorstub{audit-log}{скриншот страницы журнала аудита}
\caption{Страница журнала аудита}
\label{fig:auditlog}
\end{figure}

В отдельной вкладке доступен журнал попыток входа в систему с указанием пользователя, времени, исхода и сетевого адреса источника попытки.
Внешний вид вкладки журнала входов представлен на рисунке~\ref{fig:authlog}.

\begin{figure}[htbp]
\centering
\figureorstub{auth-log}{скриншот вкладки журнала входов в систему}
\caption{Вкладка журнала входов в систему}
\label{fig:authlog}
\end{figure}

Для просмотра истории изменений конкретной сущности на странице её редактирования предусмотрена ссылка <<История изменений>>, открывающая отфильтрованный по идентификатору сущности журнал.
В записях вида \code{update} отображается разница состояний до и после изменения, что позволяет восстановить любое предыдущее состояние сущности.
Внешний вид окна истории изменений сущности представлен на рисунке~\ref{fig:entityhistory}.

\begin{figure}[htbp]
\centering
\figureorstub{entity-history}{скриншот окна истории изменений сущности}
\caption{Окно истории изменений сущности}
\label{fig:entityhistory}
\end{figure}

\subsection{Управление учётными записями пользователей}

Страница <<Пользователи>> доступна всем аутентифицированным пользователям для просмотра, но операции создания, изменения, активации и удаления учётных записей доступны только пользователю с ролью \code{SuperAdmin}.
Внешний вид страницы пользователей представлен на рисунке~\ref{fig:users}.

\begin{figure}[htbp]
\centering
\figureorstub{users}{скриншот страницы пользователей}
\caption{Страница управления учётными записями пользователей}
\label{fig:users}
\end{figure}

Создание учётной записи выполняется через кнопку <<Добавить>>: указываются имя пользователя, начальный пароль и роль из числа предусмотренных в системе.

В строке каждой учётной записи доступны кнопки изменения, переключения признака активности и удаления.
Деактивированная учётная запись сохраняется в базе данных, но не может войти в систему.
Удаление активного пользователя с ролью \code{SuperAdmin} разрешено только при условии, что в системе остаётся хотя бы один другой активный пользователь с этой ролью~--~это исключает возможность случайной потери единственной учётной записи с правами управления пользователями.
